% from ega4-2-fr.tex
\oldpage[IV-2]{76}

\item[(ii)] Si $X$ est localement de type fini sur $k$, ou si $k'$ est une
extension de type fini de $k$, les nombres
$\operatorname{long}(\mathcal O_{x'_\alpha})$ sont des puissances de $p$.

\item[(iii)] On suppose que $X$ est localement de type fini sur $k$. Alors,
si $k'$ est parfait, tous les nombres
$\operatorname{long}(\mathcal O_{x'_\alpha})$ sont égaux et ne dépendent que
de l'extension $K$ de $k$ (et non de l'extension parfaite particulière $k'$
de $k$ considérée). En outre, il existe une extension radicielle finie $k_1$
de $k$ telle que, si $x_1$ est le point générique de
$X_1=X\otimes_k k_1$ (qui est irréductible), on ait
\[
\operatorname{long}(\mathcal O_{x'_\alpha})
=\operatorname{long}(\mathcal O_{x_1})
\quad\text{pour tout }\alpha.
\]
\end{enumerate}
\end{proposition}

\begin{proof}
\begin{enumerate}
\item[(i)] Rappelons (4.2.4) que les composantes irréductibles $X'_\alpha$,
qui sont en nombre fini, correspondent biunivoquement aux idéaux premiers
minimaux $\mathfrak p_\alpha$ de $K\otimes_k k'$, et que
$\mathcal O_{x'_\alpha}$ est un anneau artinien, isomorphe à
$(K\otimes_k k')_{\mathfrak p_\alpha}$; ces anneaux ne dépendent donc que des
extensions $K$ et $k'$ de $k$ (et sont par suite les mêmes ($k'$ étant fixé)
pour tous les $k$-préschémas intègres ayant même corps $K$ de fonctions
rationnelles). La formule (4.7.3.1) résulte de (4.7.2.1) appliquée à l'anneau
noethérien $A$ d'un voisinage affine de $x'_\alpha$ dans $X'$, et à l'anneau
$B$ d'un voisinage affine assez petit de $x''_\beta$ dans $X''$, $\mathfrak p$
et $\mathfrak q$ étant les idéaux premiers minimaux correspondant à
$x'_\alpha$ et $x''_\beta$ respectivement, et en prenant $M=A$: on notera que
$A/\mathfrak p$ est l'anneau d'un voisinage affine de $x'_\alpha$ dans
$X'_{\mathrm{red}}$, et que $B=A\otimes_{k'}k''$; si
$\mathfrak q'=\mathfrak q/\mathfrak pB$, idéal premier minimal de
$B/\mathfrak pB$, on a
\[
\mathcal O_{Z'',x''_\beta}
=(B/\mathfrak pB)_{\mathfrak q'}
=B_{\mathfrak q}/\mathfrak pB_{\mathfrak q},
\]
d'où (4.7.3.1). La dernière assertion de (i) résulte de ce que dans ce cas
$Z''$ est réduit (4.2.4), donc
$\mathfrak pB_{\mathfrak q}=\mathfrak qB_{\mathfrak q}$.

\item[(ii)] Si $k'$ est une extension de type fini de $k$, il y a une
extension transcendante pure $k'_0$ de $k$ telle que $k'$ soit une extension
finie de $k'_0$; comme $k'_0$ est séparable sur $k$, il résulte de (i)
appliqué en remplaçant $k'$ et $k''$ par $k$ et $k'_0$ respectivement, que
l'on peut remplacer $X$ par $X\otimes_k k'_0$, autrement dit, supposer que
$k'$ est une extension finie de $k$. Il y a alors une extension finie
quasi-galoisienne $k''$ de $k$ contenant $k'$, et la formule (4.7.3.1) montre
qu'il suffit de prouver l'assertion pour $k''$. Mais $k''$ est une extension
séparable d'une extension radicielle finie $k_1$ de $k$, donc, en vertu de
(i), on est ramené dans ce cas à prouver l'assertion lorsque $k'=k_1$ est en
outre radicielle. On sait alors que $K\otimes_k k_1$ est un anneau artinien
n'ayant qu'un seul idéal premier minimal (4.3.2); si $\overline K$ est une
clôture algébrique de $K$, la longueur de $K\otimes_k k_1$ divise celle de
l'anneau artinien $\overline K\otimes_k k_1$ en vertu de (4.7.1.1). Mais
$\overline K\otimes_k k_1$ est une $\overline K$-algèbre artinienne n'ayant
qu'un seul idéal premier (nécessairement maximal), et son quotient par cet
idéal est une extension finie de $\overline K$, donc nécessairement égale à
$\overline K$; la longueur de $\overline K\otimes_k k_1$ est donc égale à son
rang sur $\overline K$, c'est-à-dire à $[k_1:k]$, qui est une puissance de
$p$.

Supposons ensuite que $X$ soit localement de type fini sur $k$; alors, $X''$
est localement noethérien pour toute extension $k''$ de $k'$. Prenant pour
$k''$ une extension algébriquement close de $k'$, on est ramené, en vertu de
(4.7.3.1), à prouver l'assertion lorsque $k'$ est algébriquement clos. Comme
$k'$ est alors une extension séparable de la clôture algébrique $\overline k$
de $k$, il résulte encore de (i) qu'on est ramené au cas où
$k'=\overline k$. Or, on sait (en remplaçant $X$ par un ouvert affine de type
fini sur $k$) qu'il y a une extension radicielle

